\chapter{Теория}

\section{Разделение на сети по 2 узла}

\begin{figure}[H] 
    \centering
    \includegraphics[width=0.75\textwidth]{1.png}
    \caption{Начальные данные}
\end{figure}

Чтобы каждая сеть имела только 2 узла, нужно оставить в узловой части только 2 бита, тогда получится 4-2 = 2 узла на сеть.

\begin{figure}[H] 
    \centering
    \includegraphics[width=0.75\textwidth]{1.png}
    \caption{Начальные данные}
\end{figure}

\begin{figure}[H] 
    \centering
    \includegraphics[width=0.75\textwidth]{2.png}
    \caption{Перераспределение бит в сетевой и узловой части}
\end{figure}

\subsection{Первая сеть}
Адрес: 10.1.0.0 \\
Диапазон адресов: 10.1.0.1-10.1.0.2 \\

\subsection{Вторая сеть}
Адрес: 10.1.0.4 \\
Диапазон адресов: 10.1.0.5-10.1.0.6 \\

\subsection{Третья сеть}
Адрес: 10.1.0.4 \\
Диапазон адресов: 10.1.0.8-10.1.0.9 \\

\endinput